\documentclass[12pt]{article}
\usepackage[utf8]{inputenc}
\usepackage{geometry}
\geometry{letterpaper, margin=1in}
\usepackage{amsmath}
\usepackage{amssymb}
\usepackage{graphicx}
\usepackage{enumitem}
\usepackage{booktabs}
\usepackage{tikz}
\usepackage{array}

\title{CS 412 Assignment 6 Report}
\author{Donggyu Park}
\date{\today}

\begin{document}

\maketitle

% ==========================================
% Question 1
% ==========================================
\section*{1. Short-Answer Questions (24 points)}

\begin{enumerate}[label=(\alph*)]

    % ---- 1(a) ----
    \item \textbf{(3 points) In the K-Means algorithm, the initial placement of centroids has only a minimal impact on the quality of the resulting clustering. Discuss how initialization affects convergence and solution quality.}
    \\
    The claim is \textbf{false}. Initialization has a significant impact on both convergence speed and the quality of the final clustering.

    K-means converges to a \textit{local} optimum, not necessarily the global one. Because the algorithm greedily reassigns points to the nearest centroid and recomputes centers, the basin of attraction it follows is determined entirely by the starting positions. A poor initialization can cause:
    \begin{itemize}
        \item \textbf{Low-quality clusters:} Centroids placed in suboptimal positions lead to cluster boundaries that do not reflect the true structure of the data. For example, if two initial centroids are both placed in the same dense region, one cluster may end up covering too many points while another natural group goes unrepresented.
        \item \textbf{Wrong groupings:} Consider four points where 1 and 3 are close, and 2 and 4 are close. If the initial centroids happen to group 1 with 2 and 3 with 4, k-means will converge to that incorrect partition because no reassignment step will improve it further.
    \end{itemize}

    \textbf{K-means++ remedy:} To mitigate this, the k-means++ initialization selects the first centroid uniformly at random, then selects each subsequent centroid with probability proportional to the squared distance from the nearest already-chosen centroid:
    \[
        P(\text{select } x) \propto d(x,\, \text{nearest chosen centroid})^2.
    \]
    This spreads initial centroids across the data space, reducing the chance of a poor start. Note that it does \emph{not} deterministically pick the farthest point (which would be susceptible to outliers); instead it introduces distance-weighted randomness, making the method robust on average.

    In summary, initialization directly determines which local optimum k-means converges to, and k-means++ addresses this by favoring spread-out starting centroids.

    \vspace{2em}

    % ---- 1(b) ----
    \item \textbf{(3 points) What is the motivation behind using the k-medians algorithm instead of the k-means algorithm in certain situations? Is the k-medians algorithm more computationally demanding than k-means?}
    \\
    % TODO: Answer here.


    \vspace{2em}

    % ---- 1(c) ----
    \item \textbf{(3 points) A student claims that in a perfect clustering, the conditional entropy $H(Y \mid C)$ is maximized. Evaluate this claim and determine what value $H(Y \mid C)$ should take in the ideal case. Justify your answer.}
    \\
    % TODO: Answer here.


    \vspace{2em}

    % ---- 1(d) ----
    \item \textbf{(3 points) In the K-Means Clustering algorithm, why might a cluster center have no points assigned to it after the cluster assignment step? Explain how this situation can arise during the algorithm's iterations.}
    \\
    % TODO: Answer here.


    \vspace{2em}

    % ---- 1(e) ----
    \item \textbf{(4 points) What is the best clustering method (partition-based, hierarchical, or density) for clustering employees in a company based on their salaries and years of working experience? Justify your answer, assuming the data is continuous and free of significant outliers.}
    \\
    % TODO: Answer here.


    \vspace{2em}

    % ---- 1(f) ----
    \item \textbf{(4 points) Explain why a typical hierarchical clustering method may not lead to better quality clustering than a partitioning method.}
    \\
    % TODO: Answer here.


    \vspace{2em}

    % ---- 1(g) ----
    \item \textbf{(4 points) In DBSCAN, if a point $p$ is directly density-reachable from a point $q$, explain why this implies that $q$ must be a core point, but $p$ does not necessarily have to be.}

    \textit{Recall:} A point $p$ is directly density-reachable from $q$ (w.r.t.\ $\mathit{Eps}$, $\mathit{MinPts}$) if
    \[
        p \in N_{\mathit{Eps}}(q) \quad \text{and} \quad |N_{\mathit{Eps}}(q)| \geq \mathit{MinPts}.
    \]
    \\
    % TODO: Answer here.


\end{enumerate}

\newpage
% ==========================================
% Question 2
% ==========================================
\section*{2. Clustering Evaluation (26 points)}

Cluster information (n = 19, $|X|=6$, $|O|=6$, $|\Delta|=7$):
\begin{center}
\begin{tabular}{lccc|c}
\toprule
 & $X$ & $O$ & $\Delta$ & Total \\
\midrule
Cluster 1 & 5 & 0 & 2 & 7 \\
Cluster 2 & 1 & 4 & 1 & 6 \\
Cluster 3 & 0 & 2 & 4 & 6 \\
\bottomrule
\end{tabular}
\end{center}

\begin{enumerate}[label=(\alph*)]

    % ---- 2(a) ----
    \item \textbf{(4 points) What is the total purity of this clustering?}

    The purity is defined as:
    \[
        \text{Purity}(\Omega, \mathbb{C}) = \frac{1}{n} \sum_{k} \max_{j} |C_k \cap L_j|
    \]
    % TODO: Compute total purity here.


    \vspace{2em}

    % ---- 2(b) ----
    \item \textbf{(6 points) What is the precision of each cluster?}

    For each cluster $C_k$ with majority class $L_j$, precision is defined as:
    \[
        \text{Precision}(C_k) = \frac{|C_k \cap L_j|}{|C_k|}
    \]
    % TODO: Compute precision for each cluster.


    \vspace{2em}

    % ---- 2(c) ----
    \item \textbf{(6 points) What is the recall of each cluster for its majority class?}

    For each cluster $C_k$ with majority class $L_j$, recall is defined as:
    \[
        \text{Recall}(C_k) = \frac{|C_k \cap L_j|}{|L_j|}
    \]
    % TODO: Compute recall for each cluster.


    \vspace{2em}

    % ---- 2(d) ----
    \item \textbf{(6 points) What is the F-measure of each cluster?}

    The F-measure is defined as:
    \[
        F(C_k) = \frac{2 \cdot \text{Precision}(C_k) \cdot \text{Recall}(C_k)}{\text{Precision}(C_k) + \text{Recall}(C_k)}
    \]
    % TODO: Compute F-measure for each cluster.


    \vspace{2em}

    % ---- 2(e) ----
    \item \textbf{(4 points) What is the NMI score of this clustering? (Use natural logs.)}

    The Normalized Mutual Information is defined as:
    \[
        \text{NMI}(\Omega, \mathbb{C}) = \frac{I(\Omega; \mathbb{C})}{\frac{1}{2}[H(\Omega) + H(\mathbb{C})]}
    \]
    where $I(\Omega;\mathbb{C}) = H(\Omega) - H(\Omega \mid \mathbb{C})$.
    % TODO: Compute NMI here.


\end{enumerate}

\newpage
% ==========================================
% Question 3
% ==========================================
\section*{3. Clustering Algorithms (30 points)}

The dataset consists of 13 points:
\begin{center}
\begin{tabular}{ccc}
\toprule
Index & $x_1$ & $x_2$ \\
\midrule
1  & 1 & 3 \\
2  & 1 & 2 \\
3  & 2 & 1 \\
4  & 2 & 2 \\
5  & 2 & 3 \\
6  & 3 & 2 \\
7  & 5 & 3 \\
8  & 4 & 3 \\
9  & 4 & 5 \\
10 & 5 & 4 \\
11 & 5 & 5 \\
12 & 6 & 4 \\
13 & 6 & 5 \\
\bottomrule
\end{tabular}
\end{center}

\begin{enumerate}[label=(\alph*)]

    % ---- 3(a) ----
    \item \textbf{(10 points) Cluster the above points using the K-means algorithm with $k=2$. Use $(0,4)$ and $(6,5)$ as the initial centers. Show your reasoning.}

    \textbf{Initial centers:} $\mu_1 = (0, 4)$, $\mu_2 = (6, 5)$.

    \textbf{Iteration 1 --- Assignment step:}

    For each point, compute the Euclidean distance to each center and assign to the nearest.
    \[
        d(p, \mu) = \sqrt{(x_1 - \mu_{x_1})^2 + (x_2 - \mu_{x_2})^2}
    \]
    % TODO: Compute distances and show assignment table.


    \textbf{Iteration 1 --- Update step:}
    % TODO: Compute new centroids.


    \textbf{Iteration 2 (and subsequent):}
    % TODO: Continue until convergence. Show final clusters.


    \vspace{2em}

    % ---- 3(b) ----
    \item \textbf{(10 points) Use DBSCAN with $\mathit{MinPts} = 2$ and $\mathit{Eps} = 1.5$. Outline your clustering process.}

    \textbf{Step 1 --- Identify core points:}

    A point $p$ is a core point if $|N_{\mathit{Eps}}(p)| \geq \mathit{MinPts}$.
    % TODO: For each point, list neighbors within Eps=1.5 and identify core/border/noise.


    \textbf{Step 2 --- Form clusters:}
    % TODO: Expand clusters from core points. List final clusters.


    \vspace{2em}

    % ---- 3(c) ----
    \item \textbf{(10 points) Perform AGNES (single-link, Euclidean distance) on the above points.}

    \textbf{Single-link distance:}
    \[
        d_{\text{single}}(C_i, C_j) = \min_{p \in C_i,\, q \in C_j} d(p, q)
    \]

    \textbf{Initial state:} Each of the 13 points forms its own singleton cluster.

    % TODO: Show the merge sequence step by step.
    % At each step: identify the pair of clusters with minimum single-link distance,
    % merge them, and update the distance matrix.


\end{enumerate}

\newpage
% ==========================================
% Question 4
% ==========================================
\section*{4. Average-Link Agglomerative Clustering (20 points)}

Pairwise distances between points $a, b, c, d, e$:
\begin{center}
\begin{tabular}{c|ccccc}
 & $a$ & $b$ & $c$ & $d$ & $e$ \\ \hline
$a$ & 0     &       &       &     &   \\
$b$ & 2.1   & 0     &       &     &   \\
$c$ & 4.0   & 6.1   & 0     &     &   \\
$d$ & 7.8   & 8.6   & 6.1   & 0   &   \\
$e$ & 7.3   & 7.2   & 7.3   & 3.1 & 0 \\
\end{tabular}
\end{center}

\textbf{Average-link distance between clusters $C_i$ and $C_j$:}
\[
    d_{\text{avg}}(C_i, C_j) = \frac{1}{|C_i| \cdot |C_j|} \sum_{p \in C_i} \sum_{q \in C_j} d(p, q)
\]

\begin{enumerate}[label=(\alph*)]

    % ---- 4(a) ----
    \item \textbf{(10 points) Show which two clusters are merged at each step.}

    % TODO: Show the merge decision at each step (Steps 1 through 4).

    \textbf{Step 1:}
    % Identify the pair with minimum average-link distance and merge.


    \textbf{Step 2:}
    % Update distances and identify the next pair to merge.


    \textbf{Step 3:}
    % Update distances and identify the next pair to merge.


    \textbf{Step 4:}
    % Final merge.


    \vspace{2em}

    % ---- 4(b) ----
    \item \textbf{(10 points) Show the updated pairwise distance matrix after each merging step.}

    \textbf{After Step 1:}
    % TODO: Updated distance matrix.


    \textbf{After Step 2:}
    % TODO: Updated distance matrix.


    \textbf{After Step 3:}
    % TODO: Updated distance matrix.


\end{enumerate}

\end{document}
